% Generated draft: requires mathematical and source review.
\documentclass[11pt]{article}
\usepackage[margin=25mm]{geometry}
\usepackage{fontspec}
\setmainfont{Latin Modern Roman}
\usepackage{amsmath,amssymb,mathtools,unicode-math}
\setmathfont{Latin Modern Math}
\usepackage{enumitem,needspace,xurl,hyperref}
\setlength{\parindent}{0pt}
\setlength{\parskip}{5pt}
\setlength{\emergencystretch}{3em}
\newcommand{\N}{\mathbb{N}}
\newcommand{\Z}{\mathbb{Z}}
\newcommand{\Q}{\mathbb{Q}}
\newcommand{\R}{\mathbb{R}}
\newcommand{\C}{\mathbb{C}}
\newcommand{\F}{\mathbb{F}}
\newcommand{\A}{\mathbb{A}}
\newcommand{\PP}{\mathbb{P}}
\newcommand{\E}{\mathbb{E}}
\newcommand{\eps}{\varepsilon}
\newcommand{\dd}{\,\mathrm{d}}
\begin{document}
\section*{UnsolvedMath ID: 30005121}
\textbf{Source problem code:} catalogue entry
\subsection{Definitions and mathematical statement}
Let \(A\) be an alphabet of cardinality \(k\geq 2\), and let \(w=(w_1,\ldots,w_n)\in A^n\). A subsequence of \(w\) is a word \((w_{i_1},\ldots,w_{i_\ell})\) with \(1\leq i_1<\cdots<i_\ell\leq n\); its positions need not be consecutive. Two subsequences are \emph{twins} if they are identical as words and their sets of positions are disjoint. Thus a pair of twins of length \(\ell\) consists of indices
\[
1\leq i_1<\cdots<i_\ell\leq n,\qquad
1\leq j_1<\cdots<j_\ell\leq n,
\]
with \(\{i_1,\ldots,i_\ell\}\cap\{j_1,\ldots,j_\ell\}=\varnothing\) and \(w_{i_r}=w_{j_r}\) for every \(1\leq r\leq\ell\). Define \(t(w)\) to be the maximum such \(\ell\), allowing \(\ell=0\), and define
\[
t(n,k)=\min_{w\in A^n}t(w).
\]
The length counts the symbols in each twin separately, so \(0\leq t(w)\leq\lfloor n/2\rfloor\).

Fix \(|A|=3\), and let \(W_3(n)\) be uniformly distributed on \(A^n\). The supplied problem asks both of the following questions as \(n\to\infty\):
\begin{enumerate}
\item Is it true that, for every \(\varepsilon>0\),
\[
\lim_{n\to\infty}\mathbb{P}\!\left(t(W_3(n))\geq\left(\frac12-\varepsilon\right)n\right)=1?
\]
This expresses the supplied asymptotically almost sure assertion \(t(W_3(n))=n/2-o(n)\).
\item Is it true that
\[
\lim_{n\to\infty}\frac{t(n,3)}{n}=\frac12?
\]
Equivalently, for every \(\varepsilon>0\) there exists \(N\in\N\) such that, for every integer \(n\geq N\) and every \(w\in A^n\),
\[
t(w)\geq\left(\frac12-\varepsilon\right)n.
\]
\end{enumerate}
The extracted original uses the term ``subwords'' and contains corrupted notation. The subsequence interpretation and the uniform random model above follow the supplied clean statement; the original extraction alone does not fully specify the random model. These are questions, with the open status supplied by the record rather than independently verified here.
\subsection{Short English statement}
Must every sufficiently long ternary word contain two identical subsequences, using disjoint positions, each of length \(n/2-o(n)\)? Does the same conclusion hold with probability tending to one for a uniformly random ternary word? Length refers to each of the two subsequences. Both assertions are posed as questions requiring mathematical review.
\subsection{Sources}
\begin{itemize}
\item Supplied local UnsolvedMath registry record, permanent ID \(30005121\), title ``Long Twin Subsequences in Ternary Words,'' source problem code \texttt{catalogue entry}, in the Oberwolfach Reports Open Problems collection.
\item Supplied report citation: \emph{Combinatorics, Probability and Computing (2023)}, Oberwolfach Reports, pp.\ 1165--1237, DOI \texttt{10.4171/owr/2022/22}. \url{https://doi.org/10.4171/owr/2022/22}. The extracted original attributes the question to Andrzej Ruci\'nski; the registry's proposed year is 2022. The supplied report designation and year are retained without bibliographic reconciliation.
\item Supplied literature reference: Axenovich, Person, and Puzynina, ``twins in words.'' \url{https://doi.org/10.1016/j.jcta.2013.01.001}.
\item Supplied literature reference: Dudek, Grytczuk, and Rucinski, ``long twins in random words.'' \url{https://arxiv.org/abs/2112.14197}.
\item Literature assessment dated exactly \texttt{2026-08-22}: the supplied curation labels the problem open and states that known deterministic and random lower bounds for ternary words are strictly below \(n/2\), while the asserted \(n/2-o(n)\) behavior has not been proved. This is supplied dated curation, not an independently verified literature assessment; the record requests independent verification of time-sensitive status and cited claims.
\item The supplied registry labels the statement \texttt{corrected\_verified} and reports that the official Oberwolfach report was checked directly when preparing its clean statement. This is a source-provided verification claim; no independent source consultation is asserted in this draft.
\end{itemize}
\end{document}
